\section{Gauss–Bonnet 定理}\label{sec:gauss-bonnet-theorem}

现在利用 Cartan 结构方程证明 \textbf{Gauss–Bonnet 定理}, 将紧曲面的曲率与其拓扑联系起来. 设 $M$ 是紧定向 2 维流形, $X$ 是 $M$ 上的向量场.

\begin{RiemannDefinition}\label{def:isolated-singularity-of-vector-field}
若 $X_p=0$, 则称点 $p\in M$ 为 $X$ 的\textbf{奇点}. 若存在 $p$ 的邻域 $V\subset M$, 使得 $p$ 是 $X$ 在 $V$ 中的唯一奇点, 则称这个奇点为\textbf{孤立奇点}.
\end{RiemannDefinition}

由于 $M$ 紧, 若 $X$ 的所有奇点都孤立, 则奇点的数目有限 (否则它们将聚集于某个非孤立奇点).

可按如下方法给 $X\in\mathfrak{X}(M)$ 的每个孤立奇点 $p\in V$ 赋予一个整数, 称为 $X$ 在 $p$ 处的\textbf{指标}:
\begin{enumerate}[label=(\roman*),leftmargin=*]
\item 在 $M$ 上固定一个 Riemann 度量;
\item 选取定义在 $V\setminus\{p\}$ 上的正定向标准正交标架 $\{F_1,F_2\}$, 使得
\[
  F_1=\frac{X}{\lVert X\rVert},
\]
设 $\{\overline\omega{}^1,\overline\omega{}^2\}$ 为其对偶余标架, $\overline\omega{}_1^2$ 为相应的联络形式;
\item 必要时缩小 $V$, 取定义在 $V$ 上的正定向标准正交标架 $\{E_1,E_2\}$, 其对偶余标架为 $\{\omega^1,\omega^2\}$, 联络形式为 $\omega_1^2$;
\item 取 $p$ 在 $V$ 中同胚于圆盘且有光滑边界 $\partial D$ 的邻域 $D$, 给边界赋予诱导定向, 定义 $X$ 在 $p$ 处的指标 $I_p$ 为
\[
  2\pi I_p=\int_{\partial D}\sigma,
\]
其中 $\sigma:=\overline\omega{}_1^2-\omega_1^2$ 是 Proposition~\ref{prop:connection-form-under-oriented-frame-change} 中的形式.
\end{enumerate}
回顾 $\sigma$ 满足 $\sigma=\mathrm{d}\theta$, 其中 $\theta$ 是 $E_1$ 与 $F_1$ 的夹角. 因此 $I_p$ 必是整数. 直观而言, 向量场 $X$ 的指标衡量逆时针绕奇点一周时 $X$ 旋转的圈数, 若 $X$ 本身逆时针旋转则计为正, 否则计为负.

\begin{RiemannExample}\label{ex:four-vector-field-indices}
在 $M=\mathbb{R}^2$ 中, 下列向量场在原点有孤立奇点, 其指标如下 (参见 Fig.~\ref{fig:vector-field-singularity-indices}):
\begin{enumerate}[label=(\arabic*),leftmargin=*]
\item $X_{(x,y)}=(x,y)$ 的指标为 $1$;
\item $Y_{(x,y)}=(-y,x)$ 的指标为 $1$;
\item $Z_{(x,y)}=(y,x)$ 的指标为 $-1$;
\item $W_{(x,y)}=(x,-y)$ 的指标为 $-1$.
\end{enumerate}
\end{RiemannExample}
\RGfigure[.75\textwidth]{vector-field-singularity-indices}{计算向量场 $X$, $Y$, $Z$ 和 $W$ 的指标}

现在验证指标是良定义的. 首先注意, 由于 $\sigma$ 是闭的, $I_p$ 不依赖于 $D$ 的选取. 事实上, 任意两个这样的圆盘的边界必同伦 [参见第~\ref{chap:differential-forms}~章的 Exercise~\ref{ex:stokes-theorem}(\ref{ex:homotopy-invariance-of-integral})]. 接下来证明 $I_p$ 不依赖于标架 $\{E_1,E_2\}$ 的选取. 更确切地, 将证明
\[
  I_p=\lim_{r\to0}\frac1{2\pi}\int_{S_r(p)}\overline\omega{}_1^2,
\]
其中 $S_r(p)$ 是以 $p$ 为中心, 半径为 $r$ 的正规球面. 事实上, 若 $r_1>r_2>0$ 为正规球面的半径, 则
\begin{equation}\label{eq:annulus-curvature-and-connection}
\begin{aligned}
 &\int_{S_{r_1}(p)}\overline\omega{}_1^2-\int_{S_{r_2}(p)}\overline\omega{}_1^2
  =\int_{\Delta_{12}}\mathrm{d}\overline\omega{}_1^2\\
 &\qquad=-\int_{\Delta_{12}}K\,\overline\omega{}^1\wedge\overline\omega{}^2
  =-\int_{\Delta_{12}}K,
\end{aligned}
\end{equation}
其中 $\Delta_{12}=B_{r_1}(p)\setminus B_{r_2}(p)$. 由于 $K$ 连续, 当 $r_1\to0$ 时, 有
\[
  \left(\int_{S_{r_1}(p)}\overline\omega{}_1^2-\int_{S_{r_2}(p)}\overline\omega{}_1^2\right)\longrightarrow0.
\]
因此, 若 $\{r_n\}$ 是递减并收敛到零的正数序列, 则序列
\[
  \left\{\int_{S_{r_n}(p)}\overline\omega{}_1^2\right\}
\]
是 Cauchy 序列, 从而收敛. 令
\[
  \overline I_p:=\lim_{r\to0}\frac1{2\pi}\int_{S_r(p)}\overline\omega{}_1^2.
\]
在 \eqref{eq:annulus-curvature-and-connection} 中取 $r_2\to0$ 的极限, 得
\begin{align*}
  \int_{S_{r_1}(p)}\overline\omega{}_1^2-2\pi\overline I_p
  &=-\int_{B_{r_1}(p)}K=-\int_{B_{r_1}(p)}K\,\omega^1\wedge\omega^2\\
  &=\int_{B_{r_1}(p)}\mathrm{d}\omega_1^2=\int_{S_{r_1}(p)}\omega_1^2,
\end{align*}
因而
\[
  2\pi I_p=\int_{S_{r_1}(p)}\sigma
    =\int_{S_{r_1}(p)}\overline\omega{}_1^2-\omega_1^2=2\pi\overline I_p.
\]

最后, 证明 $I_p$ 不依赖于 Riemann 度量的选取. 事实上, 若 $\langle\cdot,\cdot\rangle_0$, $\langle\cdot,\cdot\rangle_1$ 是 $M$ 上的两个 Riemann 度量, 则容易验证
\[
  \langle\cdot,\cdot\rangle_t:=(1-t)\langle\cdot,\cdot\rangle_0+t\langle\cdot,\cdot\rangle_1
\]
也是 $M$ 上的一个 Riemann 度量, 且按度量 $\langle\cdot,\cdot\rangle_t$ 计算的指标 $I_p(t)$ 是 $t$ 的连续函数 [参见 Exercise~\ref{ex:gauss-bonnet}(\ref{ex:index-continuous-under-metric-interpolation})]. 由于对所有 $t\in[0,1]$, $I_p(t)$ 都是整数, 可知 $I_p(0)=I_p(1)$.

因此 $I_p$ 只依赖于向量场 $X\in\mathfrak{X}(M)$. 现在可以陈述 Gauss–Bonnet 定理.

\begin{RiemannTheorem}[Gauss–Bonnet]\label{thm:gauss-bonnet}
设 $M$ 是紧定向 2 维流形, $X$ 是 $M$ 中一个具有孤立奇点 $p_1,\ldots,p_k$ 的向量场. 则对 $M$ 上的任意 Riemann 度量,
\begin{equation}\label{eq:gauss-bonnet-index-formula}
  \int_M K=2\pi\sum_{i=1}^k I_{p_i}
\end{equation}
其中 $K$ 是 Gauss 曲率.
\end{RiemannTheorem}

\begin{RiemannProof}
考虑正定向标准正交标架 $\{F_1,F_2\}$, 其中
\[
  F_1=\frac{X}{\lVert X\rVert},
\]
它定义在 $M\setminus\bigcup_{i=1}^k\{p_i\}$ 上, 其对偶余标架为 $\{\overline\omega{}^1,\overline\omega{}^2\}$, 联络形式为 $\overline\omega{}_1^2$. 取充分小的 $r>0$, 令 $B_i:=B_r(p_i)$, 使得当 $i\ne j$ 时, $\overline B_i\cap\overline B_j=\varnothing$. 注意,
\begin{align*}
  \int_{M\setminus\bigcup_{i=1}^k B_i}K
  &=\int_{M\setminus\bigcup_{i=1}^k B_i}K\,\overline\omega{}^1\wedge\overline\omega{}^2
  =-\int_{M\setminus\bigcup_{i=1}^k B_i}\mathrm{d}\overline\omega{}_1^2\\
  &=\int_{\bigcup_{i=1}^k\partial B_i}\overline\omega{}_1^2
  =\sum_{i=1}^k\int_{\partial B_i}\overline\omega{}_1^2,
\end{align*}
其中 $\partial B_i$ 取由 $B_i$ 的定向诱导的定向. 取 $r\to0$ 的极限, 得
\[
  \int_M K=2\pi\sum_{i=1}^k I_{p_i}.
\]
\end{RiemannProof}

\begin{RiemannRemark}\label{rem:euler-characteristic-from-indices-and-triangulation}
\begin{enumerate}[label=(\arabic*),leftmargin=*]
\item 由于 \eqref{eq:gauss-bonnet-index-formula} 的右端与度量无关, 可知对 $M$ 上的\textbf{所有} Riemann 度量, $\int_M K$ 都相同.
\item 由于 \eqref{eq:gauss-bonnet-index-formula} 的左端与向量场 $X$ 无关, 可知对 $M$ 上所有具有孤立奇点的向量场, $\chi(M):=\sum_{i=1}^k I_{p_i}$ 都相同. 这称为 $M$ 的 \textbf{Euler 特征数}.
\item 回顾, $M$ 的一个\textbf{三角剖分}是将 $M$ 分解为有限个三角形 (即 Euclidean 三角形在参数化下的像), 使得任意两个三角形的交集或者是一条公共边, 或者是一个公共顶点, 或者为空 (可以证明这样的三角剖分总存在). 给定一个三角剖分, 可构造具有以下性质的向量场 $X$ (参见 Fig.~\ref{fig:vector-field-on-triangulation}):
\begin{enumerate}[label=(\alph*),leftmargin=*]
\item 每个顶点都是一个汇点型奇点, 即对以该奇点为中心的某些局部坐标 $(x,y)$,
\[
  X=-x\frac{\partial}{\partial x}-y\frac{\partial}{\partial y};
\]
\item 每个 2 维面的内部恰有一个源点型奇点, 即对以该奇点为中心的某些局部坐标 $(x,y)$,
\[
  X=x\frac{\partial}{\partial x}+y\frac{\partial}{\partial y};
\]
\item 每条边由向量场的积分曲线构成, 并恰好包含一个不是顶点的奇点.
\end{enumerate}
容易看出, 所有奇点都是孤立的, 顶点和 2 维面内的奇点指标为 $1$, 边上的奇点指标为 $-1$. 所以
\[
  \chi(M)=V-E+F,
\]
其中 $V$ 是任意三角剖分的顶点数, $E$ 是边数, $F$ 是 2 维面的数目. 这就是第~\ref{chap:differentiable-manifolds}~章的 Exercise~\ref{ex:topological-manifolds}(\ref{ex:euler-characteristic-of-surfaces}) 中使用的定义.
\end{enumerate}
\end{RiemannRemark}
\RGfigure[.7\textwidth]{vector-field-on-triangulation}{与三角剖分相联系的向量场}

\begin{RiemannExample}\label{ex:euler-characteristics-of-sphere-and-torus}
\begin{enumerate}[label=(\arabic*),leftmargin=*]
\item 在 $S^2$ 上选取标准度量, 得
\[
  \chi(S^2)=\frac1{2\pi}\int_{S^2}1=\frac1{2\pi}\operatorname{vol}(S^2)=2.
\]
由此可以得出一些结论:
\begin{enumerate}[label=(\alph*),leftmargin=*]
\item $S^2$ 上不存在零曲率度量, 否则将有 $\chi(S^2)=0$;
\item $S^2$ 上不存在无奇点的向量场, 因为这同样会蕴含 $\chi(S^2)=0$;
\item 对 $S^2$ 的任意三角剖分, 有 $V-E+F=2$. 特别地, 这证明了具有三角形 2 维面的凸多面体的 Euler 公式, 因为这些多面体显然给出 $S^2$ 的三角剖分.
\end{enumerate}
\item 如将在第~\RGref{sec:manifolds-of-constant-curvature}~节所见, 环面 $T^2$ 有一个零曲率度量, 因而 $\chi(T^2)=0$. 这一点也可由 $T^2$ 上存在无奇点的向量场看出.
\end{enumerate}
\end{RiemannExample}

\begin{RiemannExercise}\label{ex:gauss-bonnet}
\begin{enumerate}[label=(\arabic*),ref=\arabic*,leftmargin=*]
\item\label{ex:index-continuous-under-metric-interpolation} 证明若 $\langle\cdot,\cdot\rangle_0$, $\langle\cdot,\cdot\rangle_1$ 是 $M$ 上的两个 Riemann 度量, 则
\[
  \langle\cdot,\cdot\rangle_t:=(1-t)\langle\cdot,\cdot\rangle_0+t\langle\cdot,\cdot\rangle_1
\]
也是 $M$ 上的 Riemann 度量, 且按度量 $\langle\cdot,\cdot\rangle_t$ 计算的指标 $I_p(t)$ 是 $t$ 的连续函数.
\item\label{ex:gauss-bonnet-for-nonorientable-manifolds} (\emph{不可定向流形的 Gauss–Bonnet 定理}) 设 $(M,g)$ 是紧不可定向 2 维 Riemann 流形, $\pi:\overline M\to M$ 是其可定向二重覆盖 [参见第~\ref{chap:differentiable-manifolds}~章的 Exercise~\ref{ex:orientability}(\ref{ex:orientable-double-covering})]. 证明:
\begin{enumerate}[label=(\alph*),leftmargin=*]
\item $\chi(\overline M)=2\chi(M)$;
\item $\overline K=\pi^*K$, 其中 $\overline K$ 是 $\overline M$ 上 Riemann 度量 $\overline g:=\pi^*g$ 的 Gauss 曲率;
\item $2\pi\chi(M)=\frac12\int_{\overline M}\overline K$.
\end{enumerate}
{\small (\textbf{Remark:} 即使 $M$ 不可定向, 仍可用 $\int_Mf=\frac12\int_{\overline M}\pi^*f$ 定义 $M$ 上函数 $f$ 的积分; 在这一定义下, Gauss–Bonnet 定理也适用于不可定向的 2 维 Riemann 流形).}
\item\label{ex:gauss-bonnet-for-manifolds-with-boundary} (\emph{带边界流形的 Gauss–Bonnet 定理}) 设 $M$ 是紧定向 2 维带边界流形, $X$ 是 $M$ 中与 $\partial M$ \textbf{横截}的向量场 (即对所有 $p\in\partial M$, $X_p\notin T_p\partial M$), 具有孤立奇点 $p_1,\ldots,p_k\in M\setminus\partial M$. 证明对 $M$ 上任意 Riemann 度量,
\[
  \int_M K+\int_{\partial M}k_g=2\pi\sum_{i=1}^k I_{p_i}
\]
其中 $K$ 是 $M$ 的 Gauss 曲率, $k_g$ 是 $\partial M$ 的测地曲率.
\item\label{ex:closed-geodesics-intersect-in-positive-gauss-curvature} 设 $(M,g)$ 是具有正 Gauss 曲率的紧可定向 2 维 Riemann 流形. 证明任意两条无自交的闭测地线必相交.
\item\label{ex:hessian-and-morse-theorem} 设 $M$ 是可微流形, $f:M\to\mathbb{R}$ 是光滑函数.
\begin{enumerate}[label=(\alph*),leftmargin=*]
\item (\emph{Hessian}) 设 $p\in M$ 是 $f$ 的一个临界点 (即 $(\mathrm{d}f)_p=0$). $f$ 在 $p$ 处的 \textbf{Hessian} 是由
\[
  (Hf)_p(v,w)=\left.\frac{\partial^2}{\partial t\,\partial s}\right|_{s=t=0}(f\circ\gamma)(s,t),
\]
给出的映射 $(Hf)_p:T_pM\times T_pM\to\mathbb{R}$, 其中 $\gamma:U\subset\mathbb{R}^2\to M$ 满足 $\gamma(0,0)=p$, $\frac{\partial\gamma}{\partial s}(0,0)=v$ 和 $\frac{\partial\gamma}{\partial t}(0,0)=w$. 证明 $(Hf)_p$ 是一个良定义的对称 2-张量.
\item (\emph{Morse 定理}) 若 $(Hf)_p$ 非退化, 则称 $p$ 为\textbf{非退化临界点}. 假设 $M$ 紧且 $f$ 是一个 \textbf{Morse 函数}, 即它的所有临界点都非退化. 证明临界点只有有限个. 进一步证明, 若 $M$ 为 2 维, 则
\[
  \chi(M)=m-s+n,
\]
其中 $m$, $n$ 和 $s$ 分别为极大值点, 极小值点和鞍点的数目.
{\small (\textbf{提示:} 在 $M$ 上选取一个 Riemann 度量, 并考虑向量场 $X:=\operatorname{grad}f$).}
\end{enumerate}
\item\label{ex:gauss-curvature-by-geodesic-triangle} 设 $(M,g)$ 是 2 维 Riemann 流形, $\Delta\subset M$ 是一个\textbf{测地三角形}, 即同胚于 Euclidean 三角形且各边是测地弧的像的开集. 设 $\alpha,\beta,\gamma$ 是 $\Delta$ 的内角, 即测地线在 $\partial\Delta$ 上交点处的夹角. 对充分小的 $\Delta$, 证明
\[
  \alpha+\beta+\gamma=\pi+\int_\Delta K,
\]
其中 $K$ 是 $M$ 的 Gauss 曲率. 分别利用:
\begin{enumerate}[label=(\alph*),leftmargin=*]
\item $\int_\Delta K$ 等于向量沿 $\partial\Delta$ 平行移动一周所旋转的角度这一事实;
\item 带边界流形的 Gauss–Bonnet 定理.
\end{enumerate}
{\small (\textbf{Remark:} 这个结果给出 Gauss 曲率的另一种几何解释: $K(p)=\lim_{\Delta\to p}\frac{\alpha+\beta+\gamma-\pi}{\operatorname{vol}(\Delta)}$).}
\item\label{ex:geodesics-meet-at-most-once-in-nonpositive-gauss-curvature} 设 $(M,g)$ 是具有非正 Gauss 曲率的单连通 2 维 Riemann 流形. 证明任意两条测地线至多相交于一点.
{\small (\textbf{提示:} 注意, 若两条测地线相交于多于一点, 就会存在一个\textbf{测地二边形}, 即一个同胚于圆盘且边界由两条测地弧的像组成的开集).}
\end{enumerate}
\end{RiemannExercise}
